% Glossareinträge - Beispiele für die Vorlage

% ===================================================================
% ANLEITUNG: Dies sind Beispieleinträge für das Glossar.
% Ersetzen Sie diese durch Ihre eigenen Begriffe!
%
% Zwei Arten von Einträgen:
% 1. \newglossaryentry{label}{name={}, description={}}  - Normale Begriffe
% 2. \newacronym{label}{KURZ}{Ausgeschrieben}          - Abkürzungen
% ===================================================================

% === Beispiel: Normale Glossareinträge ===

\newglossaryentry{blockchain}{
    name={Blockchain},
    description={Verteilte, unveränderliche Datenstruktur zur sicheren Speicherung von Transaktionen ohne zentrale Autorität}
}

% === Beispiel: Akronyme ===

\newacronym{egov}{E-Gov}{E-Government}
\newglossaryentry{egov-entry}{
    name={E-Government},
    description={Digitale Abwicklung von Verwaltungsleistungen und -prozessen über elektronische Kommunikationswege}
}

\newacronym{desi}{DESI}{Digital Economy and Society Index}
\newglossaryentry{desi-entry}{
    name={Digital Economy and Society Index},
    description={Jährlicher Bericht der Europäischen Kommission zur Messung des digitalen Fortschritts der EU-Mitgliedstaaten}
}

\newacronym{json}{JSON}{JavaScript Object Notation}
\newglossaryentry{json-entry}{
    name={JavaScript Object Notation},
    description={Leichtgewichtiges Datenformat für den strukturierten Austausch von Informationen}
}


% ===================================================================
% VERWENDUNG IM TEXT:
% - Erste Verwendung: \gls{api} → schreibt "Application Programming Interface (API)" aus
% - Weitere Verwendungen: \gls{api} → zeigt nur "API"
% - Plural: \glspl{api} → "APIs"
% - Großgeschrieben: \Gls{api} → am Satzanfang
%
% Für normale Glossareinträge:
% - \gls{digitalisierung} → "Digitalisierung"
% ===================================================================